\documentclass[11pt]{article}
\usepackage[utf8]{inputenc}
\usepackage[T1]{fontenc}
\usepackage[margin=1in]{geometry}
\usepackage{amsmath,amssymb,amsthm}
\usepackage{hyperref}
\usepackage{enumitem}
\usepackage{setspace}
\usepackage{graphicx}
\usepackage{booktabs}
\usepackage[round]{natbib}

\onehalfspacing

\newtheorem{definition}{Definition}
\newtheorem{assumption}{Assumption}
\newtheorem{lemma}{Lemma}
\newtheorem{proposition}{Proposition}
\newtheorem{theorem}{Theorem}
\newtheorem{corollary}{Corollary}
\newtheorem{remark}{Remark}

\title{Formal Scientific Article on Arithmetic Deconstruction and the Amador-Kähler Coherence Framework}
\author{Research Working Draft}
\date{\today}

\begin{document}

\maketitle

\begin{abstract}
This paper presents a formal academic framework for the study of arithmetic decomposition and a generalized Amador-Kähler coherence index. The objective is to construct a mathematically disciplined model that allows for rigorous symbolic reasoning, reproducible numeric evaluation, and conceptual extension across scientific domains. We define a decomposition measure based on quadratic energy and total magnitude, and construct a bounded coherence functional that remains interpretable under uniform regimes. The results establish positivity, boundedness, and stability conditions for admissible positive sequences. The framework is designed to support both theoretical analysis and computational verification in a globally standardized research context.
\end{abstract}

\noindent\textbf{Keywords:} arithmetic decomposition, Amador-Kähler framework, coherence index, mathematical formalization, scientific methodology, reproducible proof.

\section{Introduction}
The study of structure, measure, and formal representation lies at the heart of scientific reasoning. In particular, arithmetic operations provide the elementary substrate from which empirical, formal, and computational systems are interpreted. Yet, the deeper relevance of arithmetic is not limited to the mechanics of numerical calculation. It also concerns the conditions under which representation becomes meaningful, stable, and explanatory.

This paper introduces a formal model for arithmetic decomposition as a measure of internal energetic concentration and combines it with a generalized Amador-Kähler coherence functional. The proposed framework is intentionally simple, yet sufficiently rich to support rigorous statements about positivity, boundedness, and uniform-regime behavior. It offers both a conceptual and computational basis for scientific inquiry in domains where formal structure must be balanced against interpretability.

The methodology follows a standard formal style acceptable in academic research: definitions are stated precisely, propositions are derived from explicit assumptions, and proofs are provided in a reproducible manner. The aim is to generate a research artifact that can be read, tested, and extended by other scholars without ambiguity.

\section{Formal Foundation}
Let $x=(x_1,\ldots,x_n)$ be an admissible vector in $\mathbb{R}_{>0}^n$. Define the vector mass and quadratic energy by
\begin{equation}
T(x)=\sum_{i=1}^{n}x_i, \qquad E(x)=\sum_{i=1}^{n}x_i^2.
\end{equation}

\begin{definition}
The arithmetic decomposition index of $x$ is
\begin{equation}
\mathcal{D}(x)=\frac{E(x)}{T(x)+1}.
\end{equation}
\end{definition}

The additive term $+1$ prevents degeneracy in the low-mass regime and preserves regularity when the total mass is close to zero. This is especially relevant in formalized scientific models, where exact zero conditions may otherwise lead to undefined or artificially unstable interpretations.

\begin{definition}
The Amador-Kähler coherence index is defined by
\begin{equation}
\mathcal{K}(x)=1-\exp\left(-\frac{\mathcal{D}(x)}{\max\{1,n\}}\right).
\end{equation}
\end{definition}

This index maps the decomposition measure into a bounded interval, allowing it to be interpreted as a coherence score between 0 and 1.

\section{Methodological Principles}
The construction rests on four principles that align with standard scientific practice.

\begin{enumerate}[label=\textbf{P\arabic*.}]
\item \textbf{Positivity.} All admissible magnitudes are strictly positive, ensuring that the model excludes degenerate zero-mass states unless explicitly introduced.
\item \textbf{Boundedness.} The coherence measure is normalized into the unit interval; thus, it permits comparative analysis across heterogeneous data sets.
\item \textbf{Stability.} The model remains regular under repeated uniform values, a property needed for reproducible simulation and validation.
\item \textbf{Interoperability.} The formalism can be implemented and tested in multiple languages and scientific environments without altering the underlying mathematical substance.
\end{enumerate}

\section{Core Theoretical Results}
\begin{lemma}
For every admissible vector $x\in\mathbb{R}_{>0}^n$, we have $T(x)>0$ and $E(x)\ge 0$.
\end{lemma}

\begin{proof}
Since each component satisfies $x_i>0$, their sum is strictly positive and their squared values are non-negative. Therefore $T(x)>0$ and $E(x)=\sum_i x_i^2\ge 0$.
\end{proof}

\begin{theorem}
For every admissible vector $x\in\mathbb{R}_{>0}^n$, the arithmetic decomposition index satisfies
\begin{equation}
\mathcal{D}(x)\ge 0.
\end{equation}
\end{theorem}

\begin{proof}
From Lemma 1, $E(x)\ge 0$ and $T(x)+1>0$. Therefore,
\begin{equation}
\mathcal{D}(x)=\frac{E(x)}{T(x)+1}\ge 0.
\end{equation}
This completes the proof.
\end{proof}

\begin{theorem}
For every admissible vector $x\in\mathbb{R}_{>0}^n$, the Amador-Kähler coherence index satisfies
\begin{equation}
0\le \mathcal{K}(x) < 1.
\end{equation}
\end{theorem}

\begin{proof}
Set
\begin{equation}
\alpha = \frac{\mathcal{D}(x)}{\max\{1,n\}}.
\end{equation}
By Theorem 1, $\alpha\ge 0$. Hence $e^{-\alpha}\in (0,1]$. Consequently,
\begin{equation}
\mathcal{K}(x)=1-e^{-\alpha}.
\end{equation}
Since $e^{-\alpha}>0$, we have $\mathcal{K}(x)<1$. Since $e^{-\alpha}\le 1$, we also have $\mathcal{K}(x)\ge 0$. Therefore $0\le\mathcal{K}(x)<1$.
\end{proof}

\begin{proposition}
If $x=(c,\ldots,c)$ with $n$ identical positive entries, then
\begin{equation}
\mathcal{D}(x)=\frac{n c^2}{n c + 1}.
\end{equation}
Moreover, this quantity is strictly increasing in $c$ for every fixed $n\ge 1$.
\end{proposition}

\begin{proof}
If all components equal $c$, then
\begin{equation}
E(x)=\sum_{i=1}^{n}c^2 = n c^2,
\qquad
T(x)=\sum_{i=1}^{n}c = n c.
\end{equation}
Hence,
\begin{equation}
\mathcal{D}(x)=\frac{n c^2}{n c +1}.
\end{equation}
Differentiating with respect to $c$ yields
\begin{equation}
\frac{d}{dc}\mathcal{D}(x)
=\frac{n c (n c + 2)}{(n c + 1)^2}>0,
\end{equation}
which proves strict monotonicity.
\end{proof}

\begin{corollary}
In the uniform regime, the coherence index $\mathcal{K}(x)$ is also monotone increasing in $c$.
\end{corollary}

\begin{proof}
Since $\mathcal{K}(x)=1-e^{-\mathcal{D}(x)/\max\{1,n\}}$ and the exponential map is monotone decreasing, increasing $\mathcal{D}(x)$ increases $\mathcal{K}(x)$.
\end{proof}

\section{Scientific Interpretation}
The formalism above offers a compact model for assessing how dispersed or concentrated a positive vector is relative to its total mass. The quantity $E(x)$ acts as a quadratic energy term, while $T(x)+1$ regulates the scaling to preserve regularity. The resulting ratio captures the tendency of the system to concentrate magnitude within a few components, while the coherence index compresses this quantity into a normalized score.

This interpretation is useful in settings where one wants to compare systems of different sizes but comparable structure. The model is especially natural in formal scientific workflows that require a scalar descriptor with well-defined bounds, monotonicity, and computational reproducibility.

\section{Computational Validation}
To ensure reproducibility, the model is implemented both in Python and C++. The algorithms use the same definitions and test the following conditions:
\begin{enumerate}[label=\arabic*.]
\item $\mathcal{D}(x)\ge 0$;
\item $0\le \mathcal{K}(x)<1$;
\item uniform setting: $\mathcal{D}(c,\ldots,c)=nc^2/(nc+1)$.
\end{enumerate}

The implementation is intentionally minimal, making it easy to verify across different numerical environments. This contributes to the global standards of transparency, replicability, and comparability expected by contemporary scientific communities.

\section{Conclusion}
This article introduces a mathematically consistent and computationally traceable framework for arithmetic decomposition and Amador-Kähler coherence analysis. The definitions are clear, the propositions are rigorous, and the theoretical structure is well adapted for both symbolic and numerical validation. The resulting model offers a compact but meaningful measure for studying structured positive data under a bounded and interpretable formulation.

The framework is not presented as a closed doctrine, but as a formally grounded working architecture that can be refined, generalized, and integrated into larger research programs. It is therefore appropriate for academic development, reproducible experimentation, and international scientific dissemination.

\bibliographystyle{plainnat}
\bibliography{academic_dossier}

\end{document}
